\chapter{Pruebas}\label{cap:pruebas}

\section{Introducción}
El objetivo de este capítulo es documentar cómo se validó la solución de forma
pragmática en entorno local. Las pruebas cubren el flujo completo: ingesta,
limpieza, enriquecimiento, exportación de indicadores y visualización final en
los notebooks de Zeppelin.

\section{Estrategia de validación}
Se aplicaron tres tipos de verificación:

\begin{itemize}
	\item \textbf{Pruebas funcionales de flujo}: comprobar que cada etapa produce la
	salida esperada y habilita la siguiente.
	\item \textbf{Pruebas de consistencia de datos}: verificar reglas de limpieza y
	coherencia temporal/numérica tras los filtros.
	\item \textbf{Pruebas de presentación analítica}: confirmar que los notebooks
	ejecutan correctamente y generan gráficas, indicadores y mapas interpretables.
\end{itemize}

\section{Pruebas funcionales del notebook principal}
El notebook \texttt{Análisis Taxis Nueva York.zpln} se utilizó como prueba
integrada de extremo a extremo. Los párrafos principales registrados aparecen
con estado \texttt{FINISHED} y código de ejecución \texttt{SUCCESS}.

\begin{table}[H]
	\centering
	\caption{Pruebas funcionales sobre el notebook principal}
	\begin{tabular}{|p{1.7cm}|p{6.1cm}|p{4.1cm}|p{2.4cm}|}
		\hline
		\textbf{ID} & \textbf{Caso de prueba} & \textbf{Evidencia esperada} & \textbf{Resultado} \\
		\hline
		PF-01 & Carga de Parquet y revisión de esquema & Conteo de registros y muestra de filas & Superada \\
		\hline
		PF-02 & Limpieza de registros con reglas de dominio & Dataset limpio cacheado y recuento final & Superada \\
		\hline
		PF-03 & Join geográfico con GeoJSON & Columnas de zona origen/destino en dataset final & Superada \\
		\hline
		PF-04 & Consulta horaria de viajes/tarifa/duración & Tabla de 24 horas generada en Zeppelin & Superada \\
		\hline
		PF-05 & Comparativas por tipo de viaje y pasajeros & Tablas agregadas sin errores de ejecución & Superada \\
		\hline
		PF-06 & Balance recogida-destino por distrito & Tabla de diferencias por borough generada & Superada \\
		\hline
		PF-07 & Exportación de CSVs de indicadores por zona & 14 ficheros CSV generados en \texttt{/data} & Superada \\
		\hline
	\end{tabular}
\end{table}

\section{Pruebas funcionales del notebook de mapas}
El notebook \texttt{Análisis Taxis Nueva York Mapas.zpln} se validó tras la
ejecución del notebook principal, que genera los CSVs de indicadores necesarios.

\begin{table}[H]
	\centering
	\caption{Pruebas funcionales sobre el notebook de mapas}
	\begin{tabular}{|p{1.7cm}|p{6.1cm}|p{4.1cm}|p{2.4cm}|}
		\hline
		\textbf{ID} & \textbf{Caso de prueba} & \textbf{Evidencia esperada} & \textbf{Resultado} \\
		\hline
		PM-01 & Lectura de CSVs de indicadores & DataFrames cargados sin errores & Superada \\
		\hline
		PM-02 & Carga del GeoJSON de zonas & GeoDataFrame con geometrías válidas & Superada \\
		\hline
		PM-03 & Generación de mapas coropléticos & Mapas Folium renderizados con tooltips & Superada \\
		\hline
		PM-04 & Inspección interactiva por zona & Tooltip con nombre de zona e indicador al pasar el ratón & Superada \\
		\hline
	\end{tabular}
\end{table}

\section{Validación de resultados analíticos}
Además de validar ejecución, se comprobó que los resultados clave del notebook son
coherentes entre sí y con la lógica de negocio.

\begin{itemize}
\item En la serie horaria, el mínimo de viajes aparece a las 04:00 (12{,}175) y el
máximo a las 17:00 (205{,}245), patrón consistente con movilidad urbana real.
\item La comparación 1 de enero vs resto del mes confirma un comportamiento anómalo
de madrugada en el festivo, sin romper coherencia en el resto de franjas.
\item La tipología de viaje refleja dominio de viajes estándar
(2{,}696{,}555) frente a aeropuerto (94{,}614), con mayor distancia y tarifa en este
último caso.
\item El balance territorial (recogidas - destinos) identifica a Queens como emisor
neto (+119{,}591) y a Brooklyn como receptor neto (-81{,}809), en línea con una
dinámica centro-periferia esperable.
\item Los 14 CSVs exportados contienen datos coherentes con las consultas del notebook
y sirven como entrada verificable para los mapas.
\end{itemize}

\section{Pruebas sobre reglas de calidad del dato}
Se verificó que las reglas implementadas en Spark son coherentes con el dominio
del problema y se aplican de forma explícita. En particular:

\begin{itemize}
	\item no se admiten fechas de recogida o entrega nulas;
	\item la duración del viaje debe ser positiva;
	\item distancia e importes deben ser positivos;
	\item el orden temporal \texttt{dropoff > pickup} es obligatorio;
	\item la velocidad media debe estar dentro de umbrales plausibles;
	\item \texttt{PULocationID} y \texttt{DOLocationID} deben estar en rango
	válido (1..263).
\end{itemize}

Como resultado, el flujo elimina registros inconsistentes antes de calcular
indicadores, reduciendo el riesgo de conclusiones erróneas.

\section{Pruebas de coherencia entre notebooks}
Se verificó que la cadena de dependencia entre notebooks funciona correctamente:

\begin{itemize}
	\item El notebook principal genera los 14 CSVs en \texttt{/data}.
	\item El notebook de mapas los lee y genera visualizaciones sin errores.
	\item Los valores mostrados en los tooltips de los mapas son consistentes
	con las tablas del notebook principal.
	\item La versión ejecutada del notebook de mapas mantiene las salidas
	correctas para revisión posterior.
\end{itemize}

\section{Limitaciones de las pruebas}
Las validaciones realizadas son suficientes para el alcance docente, pero presentan
limitaciones reconocidas:

\begin{itemize}
	\item no se ejecutó una batería sistemática de estrés en varios nodos físicos;
	\item la variabilidad de hardware local puede afectar tiempos absolutos;
	\item no se incluyó automatización CI/CD de pruebas de regresión completas.
\end{itemize}

Estas limitaciones no invalidan los hallazgos, pero delimitan su interpretación.

\section{Conclusiones}
Las pruebas confirman que la solución funciona de extremo a extremo: ambos notebooks
ejecutan correctamente sus bloques principales, las reglas de limpieza son coherentes,
los CSVs exportados alimentan correctamente los mapas, y las salidas generadas permiten
justificar resultados analíticos de manera trazable. Con ello, la implementación queda
validada para su defensa académica.